\chapter{Objetivos de la Investigación}
\label{ch:objetivos_ph}

\section{Objetivo General}

Diseñar e implementar un sistema de escritorio en Python que gestione pedidos de reparto y genere rutas óptimas en un grafo vial simulado de la ciudad del Cusco, integrando y comparando cuatro paradigmas algorítmicos —Greedy, Divide y Vencerás, Programación Dinámica y Backtracking— mediante análisis de complejidad Big-O y benchmarks empíricos, con el fin de determinar el paradigma más adecuado según el tamaño del problema y las restricciones operativas.

\section{Objetivos Específicos}

\begin{itemize}
    \item Implementar desde cero los algoritmos Greedy de selección por densidad, Merge Sort (Divide y Vencerás), Knapsack 0/1 y Held-Karp (Programación Dinámica), y Backtracking con poda de calles bloqueadas, verificando su correctitud mediante pruebas unitarias con \texttt{pytest}.
    \item Derivar y documentar la complejidad temporal y espacial Big-O de cada algoritmo implementado, comparando el crecimiento teórico con los tiempos de ejecución medidos experimentalmente para entre 3 y 15 pedidos activos.
    \item Desarrollar una interfaz gráfica de usuario (GUI) con Tkinter que permita visualizar el grafo vial, gestionar pedidos y repartidores, ejecutar cada algoritmo de forma interactiva, y presentar los resultados de comparación en gráficos de rendimiento.
\end{itemize}
